% Options for packages loaded elsewhere
\PassOptionsToPackage{unicode}{hyperref}
\PassOptionsToPackage{hyphens}{url}
\documentclass[
]{article}
\usepackage{xcolor}
\usepackage[margin=1in]{geometry}
\usepackage{amsmath,amssymb}
\setcounter{secnumdepth}{-\maxdimen} % remove section numbering
\usepackage{iftex}
\ifPDFTeX
  \usepackage[T1]{fontenc}
  \usepackage[utf8]{inputenc}
  \usepackage{textcomp} % provide euro and other symbols
\else % if luatex or xetex
  \usepackage{unicode-math} % this also loads fontspec
  \defaultfontfeatures{Scale=MatchLowercase}
  \defaultfontfeatures[\rmfamily]{Ligatures=TeX,Scale=1}
\fi
\usepackage{lmodern}
\ifPDFTeX\else
  % xetex/luatex font selection
\fi
% Use upquote if available, for straight quotes in verbatim environments
\IfFileExists{upquote.sty}{\usepackage{upquote}}{}
\IfFileExists{microtype.sty}{% use microtype if available
  \usepackage[]{microtype}
  \UseMicrotypeSet[protrusion]{basicmath} % disable protrusion for tt fonts
}{}
\makeatletter
\@ifundefined{KOMAClassName}{% if non-KOMA class
  \IfFileExists{parskip.sty}{%
    \usepackage{parskip}
  }{% else
    \setlength{\parindent}{0pt}
    \setlength{\parskip}{6pt plus 2pt minus 1pt}}
}{% if KOMA class
  \KOMAoptions{parskip=half}}
\makeatother
\usepackage{color}
\usepackage{fancyvrb}
\newcommand{\VerbBar}{|}
\newcommand{\VERB}{\Verb[commandchars=\\\{\}]}
\DefineVerbatimEnvironment{Highlighting}{Verbatim}{commandchars=\\\{\}}
% Add ',fontsize=\small' for more characters per line
\usepackage{framed}
\definecolor{shadecolor}{RGB}{248,248,248}
\newenvironment{Shaded}{\begin{snugshade}}{\end{snugshade}}
\newcommand{\AlertTok}[1]{\textcolor[rgb]{0.94,0.16,0.16}{#1}}
\newcommand{\AnnotationTok}[1]{\textcolor[rgb]{0.56,0.35,0.01}{\textbf{\textit{#1}}}}
\newcommand{\AttributeTok}[1]{\textcolor[rgb]{0.13,0.29,0.53}{#1}}
\newcommand{\BaseNTok}[1]{\textcolor[rgb]{0.00,0.00,0.81}{#1}}
\newcommand{\BuiltInTok}[1]{#1}
\newcommand{\CharTok}[1]{\textcolor[rgb]{0.31,0.60,0.02}{#1}}
\newcommand{\CommentTok}[1]{\textcolor[rgb]{0.56,0.35,0.01}{\textit{#1}}}
\newcommand{\CommentVarTok}[1]{\textcolor[rgb]{0.56,0.35,0.01}{\textbf{\textit{#1}}}}
\newcommand{\ConstantTok}[1]{\textcolor[rgb]{0.56,0.35,0.01}{#1}}
\newcommand{\ControlFlowTok}[1]{\textcolor[rgb]{0.13,0.29,0.53}{\textbf{#1}}}
\newcommand{\DataTypeTok}[1]{\textcolor[rgb]{0.13,0.29,0.53}{#1}}
\newcommand{\DecValTok}[1]{\textcolor[rgb]{0.00,0.00,0.81}{#1}}
\newcommand{\DocumentationTok}[1]{\textcolor[rgb]{0.56,0.35,0.01}{\textbf{\textit{#1}}}}
\newcommand{\ErrorTok}[1]{\textcolor[rgb]{0.64,0.00,0.00}{\textbf{#1}}}
\newcommand{\ExtensionTok}[1]{#1}
\newcommand{\FloatTok}[1]{\textcolor[rgb]{0.00,0.00,0.81}{#1}}
\newcommand{\FunctionTok}[1]{\textcolor[rgb]{0.13,0.29,0.53}{\textbf{#1}}}
\newcommand{\ImportTok}[1]{#1}
\newcommand{\InformationTok}[1]{\textcolor[rgb]{0.56,0.35,0.01}{\textbf{\textit{#1}}}}
\newcommand{\KeywordTok}[1]{\textcolor[rgb]{0.13,0.29,0.53}{\textbf{#1}}}
\newcommand{\NormalTok}[1]{#1}
\newcommand{\OperatorTok}[1]{\textcolor[rgb]{0.81,0.36,0.00}{\textbf{#1}}}
\newcommand{\OtherTok}[1]{\textcolor[rgb]{0.56,0.35,0.01}{#1}}
\newcommand{\PreprocessorTok}[1]{\textcolor[rgb]{0.56,0.35,0.01}{\textit{#1}}}
\newcommand{\RegionMarkerTok}[1]{#1}
\newcommand{\SpecialCharTok}[1]{\textcolor[rgb]{0.81,0.36,0.00}{\textbf{#1}}}
\newcommand{\SpecialStringTok}[1]{\textcolor[rgb]{0.31,0.60,0.02}{#1}}
\newcommand{\StringTok}[1]{\textcolor[rgb]{0.31,0.60,0.02}{#1}}
\newcommand{\VariableTok}[1]{\textcolor[rgb]{0.00,0.00,0.00}{#1}}
\newcommand{\VerbatimStringTok}[1]{\textcolor[rgb]{0.31,0.60,0.02}{#1}}
\newcommand{\WarningTok}[1]{\textcolor[rgb]{0.56,0.35,0.01}{\textbf{\textit{#1}}}}
\usepackage{longtable,booktabs,array}
\newcounter{none} % for unnumbered tables
\usepackage{calc} % for calculating minipage widths
% Correct order of tables after \paragraph or \subparagraph
\usepackage{etoolbox}
\makeatletter
\patchcmd\longtable{\par}{\if@noskipsec\mbox{}\fi\par}{}{}
\makeatother
% Allow footnotes in longtable head/foot
\IfFileExists{footnotehyper.sty}{\usepackage{footnotehyper}}{\usepackage{footnote}}
\makesavenoteenv{longtable}
\usepackage{graphicx}
\makeatletter
\newsavebox\pandoc@box
\newcommand*\pandocbounded[1]{% scales image to fit in text height/width
  \sbox\pandoc@box{#1}%
  \Gscale@div\@tempa{\textheight}{\dimexpr\ht\pandoc@box+\dp\pandoc@box\relax}%
  \Gscale@div\@tempb{\linewidth}{\wd\pandoc@box}%
  \ifdim\@tempb\p@<\@tempa\p@\let\@tempa\@tempb\fi% select the smaller of both
  \ifdim\@tempa\p@<\p@\scalebox{\@tempa}{\usebox\pandoc@box}%
  \else\usebox{\pandoc@box}%
  \fi%
}
% Set default figure placement to htbp
\def\fps@figure{htbp}
\makeatother
\setlength{\emergencystretch}{3em} % prevent overfull lines
\providecommand{\tightlist}{%
  \setlength{\itemsep}{0pt}\setlength{\parskip}{0pt}}
\usepackage{bookmark}
\IfFileExists{xurl.sty}{\usepackage{xurl}}{} % add URL line breaks if available
\urlstyle{same}
% fallback for those not using the hyperref driver hyperxmp:
\makeatletter
\@ifundefined{xmpquote}{\newcommand{\xmpquote}[1]{#1}}{}
\makeatother
\hypersetup{
  pdftitle={mensik\_pset\_6},
  pdfauthor={Kristina Mensik},
  hidelinks,
  pdfcreator={LaTeX via pandoc}}

\title{mensik\_pset\_6}
\author{Kristina Mensik}
\date{2026-10-24}

\begin{document}
\maketitle

\subsection{Problem 1 (a)}\label{problem-1-a}

\[
E[d_i(1)]
\] The expected treatment status i would have when treatment is assigned
is 1? Or do we factor in that we know 395 out of the 1,445 total
assigned to treatment actually took the treatment, the rest didn't
(ITT\_D) is approximately 0.273.

\begin{Shaded}
\begin{Highlighting}[]
\DecValTok{395}\SpecialCharTok{/}\DecValTok{1445}
\end{Highlighting}
\end{Shaded}

\begin{verbatim}
## [1] 0.2733564
\end{verbatim}

\[
E[Y_i(0)|d_i(1) = 0]
\] The expected untreated potential outcome for never takers is 36.48

\[
E[Y_i(0)|d_i(1) = 1]
\] The expected untreated potential outcome for compliers is 37.54.

\[
E[Y_i(1)|d_i(1) = 1]
\]

The expected treated potential outcome for compliers is 54.43.

\subsection{Problem 2}\label{problem-2}

Half compliers, half never takers (not actually treated in 1SNC)

\begin{Shaded}
\begin{Highlighting}[]
\NormalTok{schedule }\OtherTok{\textless{}{-}} \FunctionTok{data.frame}\NormalTok{(}
  \AttributeTok{Subjects =} \DecValTok{1}\SpecialCharTok{:}\DecValTok{6}\NormalTok{,}
  \AttributeTok{Type =} \FunctionTok{c}\NormalTok{(}\StringTok{"Complier"}\NormalTok{,}\StringTok{"Complier"}\NormalTok{,}\StringTok{"Complier"}\NormalTok{,}\StringTok{"NT"}\NormalTok{,}\StringTok{"NT"}\NormalTok{,}\StringTok{"NT"}\NormalTok{),}
  \AttributeTok{Yi\_0 =} \FunctionTok{c}\NormalTok{(}\DecValTok{6}\NormalTok{, }\DecValTok{6}\NormalTok{, }\DecValTok{8}\NormalTok{, }\DecValTok{1}\NormalTok{, }\DecValTok{7}\NormalTok{, }\DecValTok{7}\NormalTok{),}
  \AttributeTok{Yi\_1 =} \FunctionTok{c}\NormalTok{(}\DecValTok{3}\NormalTok{, }\DecValTok{1}\NormalTok{, }\DecValTok{2}\NormalTok{, }\DecValTok{7}\NormalTok{, }\DecValTok{7}\NormalTok{, }\DecValTok{30}\NormalTok{),}
  \AttributeTok{di\_1 =} \FunctionTok{c}\NormalTok{(}\DecValTok{1}\NormalTok{,}\DecValTok{1}\NormalTok{,}\DecValTok{1}\NormalTok{,}\DecValTok{0}\NormalTok{,}\DecValTok{0}\NormalTok{,}\DecValTok{0}\NormalTok{),}
  \AttributeTok{di\_0 =} \FunctionTok{c}\NormalTok{(}\DecValTok{0}\NormalTok{,}\DecValTok{0}\NormalTok{,}\DecValTok{0}\NormalTok{,}\DecValTok{0}\NormalTok{,}\DecValTok{0}\NormalTok{,}\DecValTok{0}\NormalTok{)}
\NormalTok{)}

\FunctionTok{View}\NormalTok{(schedule)}


\NormalTok{schedule}\SpecialCharTok{$}\NormalTok{tau }\OtherTok{\textless{}{-}}\NormalTok{ schedule}\SpecialCharTok{$}\NormalTok{Yi\_1 }\SpecialCharTok{{-}}\NormalTok{ schedule}\SpecialCharTok{$}\NormalTok{Yi\_0}
\NormalTok{ATE }\OtherTok{\textless{}{-}} \FunctionTok{mean}\NormalTok{(schedule}\SpecialCharTok{$}\NormalTok{tau)}
\NormalTok{CACE }\OtherTok{\textless{}{-}} \FunctionTok{mean}\NormalTok{(schedule}\SpecialCharTok{$}\NormalTok{tau[schedule}\SpecialCharTok{$}\NormalTok{type }\SpecialCharTok{==} \StringTok{"Complier"}\NormalTok{])}
\end{Highlighting}
\end{Shaded}

The ATE is misleading because a subject assigned to treatment that is a
never taker skews the ATE positive.

The CACE might be misleading because it does not provide information
about never takers. It's possible that if never takers did take the
treatment, the effect we are estimating would be different. However, I'm
not sure if that is more artifact of the CACE as it is the trouble of
never-takers themselves.

\subsection{Problem 3}\label{problem-3}

\subsubsection{3a}\label{a}

If the ITT is negative, the CACE must be negative is true (ITT = ITT\_D
* CACE)

\subsubsection{3b}\label{b}

No - the numerator might not stay constant when the ITTD changes, so we
can't assume that the CACE would grow/shrink when ITTD shrinks/grows.

\subsubsection{3c}\label{c}

Yes - Y(d=1) effectively does not exist, or can't be observed. The setup
is ``unidentified.''

\subsection{Problem 4}\label{problem-4}

\subsubsection{4a}\label{a-1}

The expected potential outcomes of subjects when assigned to treatment
does equal the expected potential outcome among subjects assigned to
treatment, given that they are `treatment compliers' -- these are equal
quantities under excludability.

\subsubsection{4b}\label{b-1}

Yes, because both sides of the equation have d\_i(1) = 0.

\subsubsection{4c}\label{c-1}

No - these are outcomes of people who don't take the treatment when
assigned and people who do take the treatment when assigned. We can't
expect the potential outcomes of people who do not take a treatment to
be the same as those of people who do take a treatment.

\subsubsection{4d}\label{d}

Yes. The meaningful difference in these terms is that they show
different assignment status. Excludability lets us only care about
potential outcomes vis-a-vis treatments, not treatment assignments.

\subsection{Problem 5}\label{problem-5}

We can identify compliers: those assigned to treatment who take
treatment. That's because in one sided non-comp, we rule out always by
assumption.

Among those assigned to control, we can't distinguish who is a complier
versus a never taker (both who have di(0)=0 ) unless a placebo is
assigned to the control. It also seems these statuses really depend on
the experiment -- they are not about what a subject would do in all
contexts, and even very similar versions of the same experiment.

\subsection{Problem 6}\label{problem-6}

\subsubsection{6a}\label{a-2}

\begin{Shaded}
\begin{Highlighting}[]
\DecValTok{400}\SpecialCharTok{/}\DecValTok{1000}\SpecialCharTok{*}\DecValTok{100}
\end{Highlighting}
\end{Shaded}

\begin{verbatim}
## [1] 40
\end{verbatim}

40

\begin{Shaded}
\begin{Highlighting}[]
\NormalTok{(}\DecValTok{700}\SpecialCharTok{/}\DecValTok{2000}\NormalTok{) }\SpecialCharTok{*} \DecValTok{100}
\end{Highlighting}
\end{Shaded}

\begin{verbatim}
## [1] 35
\end{verbatim}

\begin{Shaded}
\begin{Highlighting}[]
\DecValTok{40{-}35}
\end{Highlighting}
\end{Shaded}

\begin{verbatim}
## [1] 5
\end{verbatim}

\begin{Shaded}
\begin{Highlighting}[]
\DecValTok{5}\SpecialCharTok{/}\NormalTok{(}\DecValTok{500}\SpecialCharTok{/}\DecValTok{1000}\NormalTok{)}
\end{Highlighting}
\end{Shaded}

\begin{verbatim}
## [1] 10
\end{verbatim}

CACE = 10 percentage points

\subsubsection{6b}\label{b-2}

250 treated -- the ITT stays the same, but the ITT\_D changes or is now
250/1000 =0.25. The CACE then is 5/0.25, or:

\begin{Shaded}
\begin{Highlighting}[]
\DecValTok{5}\SpecialCharTok{/}\FloatTok{0.25}
\end{Highlighting}
\end{Shaded}

\begin{verbatim}
## [1] 20
\end{verbatim}

20 percentage points.

\subsubsection{6c}\label{c-2}

If effectiveness is defined as the CACE, then their exaggeration
diminishes the reported effect on turnout of their efforts. Given that
the ITT is the same whether the true contact rate is reported, it might
flummox researchers particularly if they have not themselves canvassed
on a particularly hot summer day or in a less than entirely
friendly-to-the-cause neighborhood.

\subsection{Problem 8}\label{problem-8}

\begin{Shaded}
\begin{Highlighting}[]
\FunctionTok{getwd}\NormalTok{()}
\end{Highlighting}
\end{Shaded}

\begin{verbatim}
## [1] "/Users/kom11/Documents/GitHub/regression-study/Experiments"
\end{verbatim}

\begin{Shaded}
\begin{Highlighting}[]
\NormalTok{q6data }\OtherTok{\textless{}{-}} \FunctionTok{read.csv}\NormalTok{(}\StringTok{"./doi{-}10.60600{-}yu{-}vuh8kd/Rind\_Bordia\_JASP\_1996.csv"}\NormalTok{)}
\end{Highlighting}
\end{Shaded}

\subsubsection{a}\label{a-3}

\begin{Shaded}
\begin{Highlighting}[]
\NormalTok{Y1\_treated }\OtherTok{\textless{}{-}}\NormalTok{ (}\DecValTok{591}\SpecialCharTok{/}\DecValTok{1849}\NormalTok{)}\SpecialCharTok{*}\DecValTok{100}

\NormalTok{Y0\_control }\OtherTok{\textless{}{-}}\NormalTok{ (}\DecValTok{377}\SpecialCharTok{/}\DecValTok{1430}\NormalTok{)}\SpecialCharTok{*}\DecValTok{100}

\NormalTok{ITT }\OtherTok{\textless{}{-}}\NormalTok{ Y1\_treated }\SpecialCharTok{{-}}\NormalTok{ Y0\_control}

\NormalTok{ITT}
\end{Highlighting}
\end{Shaded}

\begin{verbatim}
## [1] 5.599587
\end{verbatim}

Average recycling rates amongst those assigned to treatment were 5.6
percentage points greater than those assigned to control (not canvassed
about recycling).

\subsubsection{b}\label{b-3}

\begin{Shaded}
\begin{Highlighting}[]
\NormalTok{ITT\_D }\OtherTok{\textless{}{-}} \DecValTok{1015} \SpecialCharTok{/} \DecValTok{1849}
\NormalTok{ITT\_D}
\end{Highlighting}
\end{Shaded}

\begin{verbatim}
## [1] 0.5489454
\end{verbatim}

Among those assigned to treatment, the compliance rate is 54.89.

\subsubsection{c}\label{c-3}

{\def\LTcaptype{none} % do not increment counter
\begin{longtable}[]{@{}lll@{}}
\toprule\noalign{}
& Assigned treatment & Assigned control \\
\midrule\noalign{}
\endhead
\bottomrule\noalign{}
\endlastfoot
Total & 1,849 & 1,430 \\
Canvassed (\$d\_i\$) & 1,015 & 0 \\
Not Canvassed & 834 & 1,430 \\
Recycled (\$Y\_i\$) & 591 & 377 \\
Recycling Rate & 591/1849 = 31.96\% & 337/1430 = 26.36\% \\
& & \\
\end{longtable}
}

Expected recycling rate among those in the control group:

\[
E[Y_i(z=0)] = \pi_CE[Y_i(d=0)|C] + (1 - \pi_c)E[Y_i(d=0)|N]
\]

Expected recycling rate for those assigned to treatment:

\[
E[Y_i(z = 1)] = \pi_c * E[Y_i(d=1)|C] + (1 - \pi_c)E[Y_i(d=0)|N]
\]

\[
\pi_c = P(d_i(1) = 1)
\]

= 0.549 (the proportion of compliers).

Never-takers treatment receipt does not change across the two
expectations. Also, we know that
\(ITT_Y = E[Y_i(z = 1)] - E[Y_i(z=0)]\). We can take out the never
takers, subtract the two expected rate equations, and substitute what's
left into the above ITTY:

\[
ITT_Y = \pi_C(E[Y_i(1)|C] - E[Y_i(0)|C] + (1-\pi_C)(E[Y_i(0)|N] - E[Y_i(0) |N])
\]

sans never takers leaves:

\$\$

\pi\_C(E{[}Y\_i(1)\textbar C{]} - E{[}Y\_i(0)\textbar C{]}

\$\$

which is the CACE times the compliance rate (5.49).

\subsubsection{d}\label{d-1}

\begin{Shaded}
\begin{Highlighting}[]
\NormalTok{CACE }\OtherTok{\textless{}{-}}\NormalTok{ ITT}\SpecialCharTok{/}\NormalTok{ITT\_D}

\NormalTok{CACE}
\end{Highlighting}
\end{Shaded}

\begin{verbatim}
## [1] 10.20063
\end{verbatim}

There is a 10.2 percentage point effect of treatment for those who
complied with it (10.2 percentage point increase in recycling after
being canvassed).

\subsubsection{e}\label{e}

Comparing the recycling rates of the treated and untreated subjects
alone tends to produce misleading estimates of the CACE because we are
not isolating effects of having received treatment.

\subsection{Problem 10}\label{problem-10}

\subsubsection{a}\label{a-4}

\begin{Shaded}
\begin{Highlighting}[]
\FunctionTok{getwd}\NormalTok{()}
\end{Highlighting}
\end{Shaded}

\begin{verbatim}
## [1] "/Users/kom11/Documents/GitHub/regression-study/Experiments"
\end{verbatim}

\begin{Shaded}
\begin{Highlighting}[]
\NormalTok{data\_10 }\OtherTok{\textless{}{-}} \FunctionTok{read.csv}\NormalTok{(}\StringTok{"/Users/kom11/Documents/GitHub/regression{-}study/Experiments/Guan\_Green\_CPS\_2006.csv"}\NormalTok{)}

\NormalTok{data\_10 }\OtherTok{\textless{}{-}} \FunctionTok{na.omit}\NormalTok{(data\_10)}

\FunctionTok{table}\NormalTok{(data\_10}\SpecialCharTok{$}\NormalTok{treat2, data\_10}\SpecialCharTok{$}\NormalTok{contact)}
\end{Highlighting}
\end{Shaded}

\begin{verbatim}
##    
##        0    1
##   0 1334    0
##   1  307 2381
\end{verbatim}

\begin{Shaded}
\begin{Highlighting}[]
\NormalTok{ITT }\OtherTok{\textless{}{-}} \FunctionTok{mean}\NormalTok{(data\_10}\SpecialCharTok{$}\NormalTok{turnout[data\_10}\SpecialCharTok{$}\NormalTok{treat2 }\SpecialCharTok{==} \DecValTok{1}\NormalTok{] }\SpecialCharTok{{-}} \FunctionTok{mean}\NormalTok{(data\_10}\SpecialCharTok{$}\NormalTok{turnout[data\_10}\SpecialCharTok{$}\NormalTok{treat2 }\SpecialCharTok{==} \DecValTok{0}\NormalTok{]))}

\NormalTok{ITT}
\end{Highlighting}
\end{Shaded}

\begin{verbatim}
## [1] 0.1319296
\end{verbatim}

The ITT is 0.1319.

\subsubsection{b}\label{b-4}

RI to test sharp null that ITT is zero for all observations, taking into
account that random assignment was clustered by dorm room Yi(Z=1) =
Yi(Z=0) for each i

\begin{Shaded}
\begin{Highlighting}[]
\CommentTok{\# data\_10 \%\textgreater{}\%}
\CommentTok{\#   group\_by(dormid) \%\textgreater{}\%}
\CommentTok{\#   mutate(}
\CommentTok{\#     dorm\_pop = n()}
\CommentTok{\#   ) \%\textgreater{}\%}
\CommentTok{\#   ungroup ()}

\CommentTok{\# I\textquotesingle{}m confused how to count treated dorms. RI by dorm means that what\textquotesingle{}s supposed to happen is in a dorm where one person is treated, everyone else should be treated. }

\NormalTok{data\_10 }\SpecialCharTok{\%\textgreater{}\%}
  \FunctionTok{group\_by}\NormalTok{(dormid) }\SpecialCharTok{\%\textgreater{}\%}
  \FunctionTok{summarise}\NormalTok{(}
    \AttributeTok{any\_contact =} \FunctionTok{any}\NormalTok{(contact }\SpecialCharTok{==} \DecValTok{1}\NormalTok{, }\AttributeTok{na.rm =} \ConstantTok{TRUE}\NormalTok{)}
\NormalTok{  ) }\SpecialCharTok{\%\textgreater{}\%}
  \FunctionTok{summarise}\NormalTok{(}
    \AttributeTok{n\_dorms\_contact =} \FunctionTok{sum}\NormalTok{(any\_contact)}
\NormalTok{  )}
\end{Highlighting}
\end{Shaded}

\begin{verbatim}
## # A tibble: 1 x 1
##   n_dorms_contact
##             <int>
## 1             590
\end{verbatim}

\begin{Shaded}
\begin{Highlighting}[]
\CommentTok{\# 590 }


\NormalTok{treated\_dorms }\OtherTok{\textless{}{-}}\NormalTok{ data\_10[data\_10}\SpecialCharTok{$}\NormalTok{treat2 }\SpecialCharTok{==} \DecValTok{1}\NormalTok{, ]}

\FunctionTok{length}\NormalTok{(}\FunctionTok{unique}\NormalTok{(treated\_dorms}\SpecialCharTok{$}\NormalTok{dormid))}
\end{Highlighting}
\end{Shaded}

\begin{verbatim}
## [1] 670
\end{verbatim}

\begin{Shaded}
\begin{Highlighting}[]
\CommentTok{\# 670 dorms assigned to treatment}

\NormalTok{dorm\_check }\OtherTok{\textless{}{-}}\NormalTok{ data\_10 }\SpecialCharTok{\%\textgreater{}\%}
  \FunctionTok{group\_by}\NormalTok{(dormid) }\SpecialCharTok{\%\textgreater{}\%}
  \FunctionTok{summarise}\NormalTok{(}
    \AttributeTok{n\_assignments =} \FunctionTok{n\_distinct}\NormalTok{(treat2),}
    \AttributeTok{assignment =} \FunctionTok{first}\NormalTok{(treat2),}
    \AttributeTok{.groups =} \StringTok{"drop"}
\NormalTok{  )}

\CommentTok{\# assignment clustered by dorm room?}
\FunctionTok{all}\NormalTok{(dorm\_check}\SpecialCharTok{$}\NormalTok{n\_assignments }\SpecialCharTok{==} \DecValTok{1}\NormalTok{)}
\end{Highlighting}
\end{Shaded}

\begin{verbatim}
## [1] TRUE
\end{verbatim}

\begin{Shaded}
\begin{Highlighting}[]
\CommentTok{\# n dorms assigned to treatment}
\NormalTok{n\_dorms\_assigned }\OtherTok{\textless{}{-}} \FunctionTok{sum}\NormalTok{(dorm\_check}\SpecialCharTok{$}\NormalTok{assignment }\SpecialCharTok{==} \DecValTok{1}\NormalTok{)}

\CommentTok{\# n dorms rooms in experiment}
\NormalTok{n\_dorms }\OtherTok{\textless{}{-}} \FunctionTok{nrow}\NormalTok{(dorm\_check)}

\FunctionTok{c}\NormalTok{(}\AttributeTok{n\_dorms =}\NormalTok{ n\_dorms, }\AttributeTok{n\_dorms\_assigned =}\NormalTok{ n\_dorms\_assigned)}
\end{Highlighting}
\end{Shaded}

\begin{verbatim}
##          n_dorms n_dorms_assigned 
##             1004              670
\end{verbatim}

\begin{Shaded}
\begin{Highlighting}[]
\CommentTok{\# 670 of 1004 dorms assigned}

\FunctionTok{set.seed}\NormalTok{(}\DecValTok{6081709}\NormalTok{)}

\NormalTok{delcare\_10 }\OtherTok{\textless{}{-}} \FunctionTok{declare\_ra}\NormalTok{(}
  \AttributeTok{N =} \FunctionTok{nrow}\NormalTok{(data\_10),}
  \AttributeTok{clusters =}\NormalTok{ data\_10}\SpecialCharTok{$}\NormalTok{dormid,}
  \AttributeTok{m =}\NormalTok{ n\_dorms\_assigned}
\NormalTok{)}

\NormalTok{delcare\_10}
\end{Highlighting}
\end{Shaded}

\begin{verbatim}
## Random assignment procedure: Cluster random assignment 
## Number of units: 4022 
## Number of clusters: 1004
## Number of treatment arms: 2 
## The possible treatment categories are 0 and 1.
## The number of possible random assignments is approximately infinite. 
## The probabilities of assignment are constant across units: 
##    prob_0    prob_1 
## 0.3326693 0.6673307
\end{verbatim}

\begin{Shaded}
\begin{Highlighting}[]
\NormalTok{permmatrix\_10 }\OtherTok{\textless{}{-}} \FunctionTok{obtain\_permutation\_matrix}\NormalTok{(delcare\_10)}

\FunctionTok{dim}\NormalTok{(permmatrix\_10)}
\end{Highlighting}
\end{Shaded}

\begin{verbatim}
## [1]  4022 10000
\end{verbatim}

\begin{Shaded}
\begin{Highlighting}[]
\NormalTok{estimates\_10 }\OtherTok{\textless{}{-}} \FunctionTok{numeric}\NormalTok{(}\FunctionTok{ncol}\NormalTok{(permmatrix\_10))}

\NormalTok{itt\_hat }\OtherTok{\textless{}{-}} \FunctionTok{mean}\NormalTok{(data\_10}\SpecialCharTok{$}\NormalTok{turnout[data\_10}\SpecialCharTok{$}\NormalTok{treat2 }\SpecialCharTok{==} \DecValTok{1}\NormalTok{] }\SpecialCharTok{{-}} \FunctionTok{mean}\NormalTok{(data\_10}\SpecialCharTok{$}\NormalTok{turnout[data\_10}\SpecialCharTok{$}\NormalTok{treat2 }\SpecialCharTok{==} \DecValTok{0}\NormalTok{]))}
  
\ControlFlowTok{for}\NormalTok{ (i }\ControlFlowTok{in} \FunctionTok{seq\_len}\NormalTok{(}\FunctionTok{ncol}\NormalTok{(permmatrix\_10))) \{}
  
\NormalTok{  ri }\OtherTok{\textless{}{-}}\NormalTok{ permmatrix\_10[, i]}
  
\NormalTok{  estimates\_10 [i] }\OtherTok{\textless{}{-}} 
    \FunctionTok{mean}\NormalTok{(data\_10}\SpecialCharTok{$}\NormalTok{turnout[ri }\SpecialCharTok{==} \DecValTok{1}\NormalTok{], }\AttributeTok{na.rm =} \ConstantTok{TRUE}\NormalTok{) }\SpecialCharTok{{-}} \FunctionTok{mean}\NormalTok{(data\_10}\SpecialCharTok{$}\NormalTok{turnout[ri }\SpecialCharTok{==}\DecValTok{0}\NormalTok{], }\AttributeTok{na.rm =} \ConstantTok{TRUE}\NormalTok{)}
\NormalTok{\}}

\FunctionTok{mean}\NormalTok{(estimates\_10)}
\end{Highlighting}
\end{Shaded}

\begin{verbatim}
## [1] -0.0002462874
\end{verbatim}

\begin{Shaded}
\begin{Highlighting}[]
\CommentTok{\# centered around 0}


\NormalTok{p\_val\_10 }\OtherTok{\textless{}{-}}\NormalTok{ (}\FunctionTok{sum}\NormalTok{(}\FunctionTok{abs}\NormalTok{(estimates\_10) }\SpecialCharTok{\textgreater{}=} \FunctionTok{abs}\NormalTok{(itt\_hat)) }\SpecialCharTok{+} \DecValTok{1}\NormalTok{)}\SpecialCharTok{/}\NormalTok{(}\FunctionTok{length}\NormalTok{(estimates\_10)}\SpecialCharTok{+}\DecValTok{1}\NormalTok{) }\CommentTok{\#at first got 0, but then used mc correction bc SEs likely dwarfed by large effect size}

\NormalTok{itt\_hat}
\end{Highlighting}
\end{Shaded}

\begin{verbatim}
## [1] 0.1319296
\end{verbatim}

\begin{Shaded}
\begin{Highlighting}[]
\NormalTok{p\_val\_10}
\end{Highlighting}
\end{Shaded}

\begin{verbatim}
## [1] 9.999e-05
\end{verbatim}

The effect of assignment to canvassing was significantly statistic, so
we can reject the null hypothesis.

\subsubsection{c}\label{c-4}

CACE assuming the leaflet alone had no effect on turnout seems like the
functional equivalent of the exclusion restriction (thing that matters
is being treated where treated means canvasser interaction). It is:

\begin{Shaded}
\begin{Highlighting}[]
\NormalTok{itt\_hat}\SpecialCharTok{/}\NormalTok{(}\FunctionTok{sum}\NormalTok{(data\_10}\SpecialCharTok{$}\NormalTok{contact)}\SpecialCharTok{/}\FunctionTok{sum}\NormalTok{(data\_10}\SpecialCharTok{$}\NormalTok{treat2))}
\end{Highlighting}
\end{Shaded}

\begin{verbatim}
## [1] 0.1489402
\end{verbatim}

\subsubsection{d}\label{d-2}

Turnout rate would have been 0.7 instead of 0.8

\begin{Shaded}
\begin{Highlighting}[]
\CommentTok{\#turnout for assigned to treatment}
\DecValTok{2152}\SpecialCharTok{/}\DecValTok{2688}
\end{Highlighting}
\end{Shaded}

\begin{verbatim}
## [1] 0.8005952
\end{verbatim}

\begin{Shaded}
\begin{Highlighting}[]
\NormalTok{treated\_turnout }\OtherTok{\textless{}{-}} \DecValTok{2152}\SpecialCharTok{/}\DecValTok{2688}

\NormalTok{no\_flyer\_treated\_turnout }\OtherTok{\textless{}{-}}\NormalTok{ (}\FloatTok{0.79} \SpecialCharTok{*} \DecValTok{2688}\NormalTok{)}

\NormalTok{no\_flyer\_treated\_turnout }
\end{Highlighting}
\end{Shaded}

\begin{verbatim}
## [1] 2123.52
\end{verbatim}

\begin{Shaded}
\begin{Highlighting}[]
\CommentTok{\#assuming the .52 means 1 percentage point is rounded}

\NormalTok{no\_flyer\_treated\_turnout }\OtherTok{\textless{}{-}}\NormalTok{ treated\_turnout }\SpecialCharTok{{-}} \FloatTok{0.01}

\NormalTok{no\_flyer\_treated\_turnout}
\end{Highlighting}
\end{Shaded}

\begin{verbatim}
## [1] 0.7905952
\end{verbatim}

\begin{Shaded}
\begin{Highlighting}[]
\CommentTok{\#this was sort of silly to do because I was just told it is 0.8{-}0.01.}

\CommentTok{\# total proportion compliers}
\DecValTok{2380}\SpecialCharTok{/}\DecValTok{2688}
\end{Highlighting}
\end{Shaded}

\begin{verbatim}
## [1] 0.8854167
\end{verbatim}

\$\$

\pi\_\{d(1)=1\} = 2380/2688

\textbackslash{} \textbackslash{}

\pi\_\{d(1)=1\} = \frac{Y(d = 1, z = 1)}{Y(z = 1)}

\textbackslash{} \textbackslash{}

\mu\_\{comlier-canvassed-human\} = E{[}Y(z=1, d=1)\textbar d(1) = 1{]}

\textbackslash{} \textbackslash{}

\mu\_\{complier-not-canvassed\} = E{[}Y(z=0, d = 0)\textbar d(1) = 1{]}

\textbackslash{} \textbackslash{}

\mu\_\{nevertaker-not-canvassed\} = E{[}Y(z=0, d = 0)\textbar d(1) =
0{]}

\$\$ Treatment-assigned turnout then:

\$\$

E{[}Y\_i\textbar z\_i = 1{]} =
\pi\emph{\{d(1)=1\}\mu}\{comlier-canvassed-human\} +
(1-\pi\emph{\{d(1)=1\})\mu}\{nevertaker-not-canvassed\} + 0.01

\$\$ Control-assigned turnout:

\$\$

E{[}Y\_i\textbar z\_i = 0{]} =
\pi\emph{\{d(1)=1\}\mu}\{complier-not-canvassed\} + (1 - \pi\emph{\{d(1)
= 1\})\mu}\{nevertaker-not-canvassed\}

\$\$ \#\#\# e CACE of canvassing

I know E{[}{]} treated = 0.7906, and E{[}{]} control = 0.6687

\begin{Shaded}
\begin{Highlighting}[]
\NormalTok{prop\_complier }\OtherTok{\textless{}{-}} \FunctionTok{mean}\NormalTok{(data\_10}\SpecialCharTok{$}\NormalTok{contact[data\_10}\SpecialCharTok{$}\NormalTok{treat2 }\SpecialCharTok{==} \DecValTok{1}\NormalTok{])}

\CommentTok{\# itt hat {-} 0.01 / pi complier}
\CommentTok{\# }
\CommentTok{\# cace\_10e \textless{}{-} itt\_hat {-} 0.01 /pi\_complier}

\NormalTok{control\_turnout }\OtherTok{\textless{}{-}} \FunctionTok{mean}\NormalTok{(data\_10}\SpecialCharTok{$}\NormalTok{turnout[data\_10}\SpecialCharTok{$}\NormalTok{treat2 }\SpecialCharTok{==} \DecValTok{0}\NormalTok{])}

\CommentTok{\#treated\_turnout }
\CommentTok{\#no\_flyer\_treated\_turnout }

\NormalTok{cace\_canvassing }\OtherTok{\textless{}{-}}\NormalTok{ (}
\NormalTok{  no\_flyer\_treated\_turnout }\SpecialCharTok{{-}}\NormalTok{ control\_turnout}
\NormalTok{) }\SpecialCharTok{/}\NormalTok{ prop\_complier}

\NormalTok{cace\_canvassing}
\end{Highlighting}
\end{Shaded}

\begin{verbatim}
## [1] 0.1376509
\end{verbatim}

The CACE is 0.137

\subsubsection{f}\label{f}

CACE if leaflet had no effect on compliers and raised turnout among
never takers by 3 pct points. Have: ITT\_hat proportion compliers
(prop\_complier) Y\_hat for control (control\_turnout) Y-hat for treated
when all assigned to treatment are counted as treated (treated\_turnout)

\begin{Shaded}
\begin{Highlighting}[]
\CommentTok{\#change to ITT is 0.03 * prop never takers}
\CommentTok{\# is 0.03 * (1{-}prop\_complier)}

\NormalTok{prop\_nts }\OtherTok{\textless{}{-}} \DecValTok{1}\SpecialCharTok{{-}}\NormalTok{prop\_complier}

\NormalTok{new\_itt }\OtherTok{\textless{}{-}}\NormalTok{ itt\_hat}\SpecialCharTok{{-}}\NormalTok{(.}\DecValTok{03}\SpecialCharTok{*}\NormalTok{prop\_nts)}

\CommentTok{\# cace is new / prop\_complier}

\NormalTok{new\_itt }\SpecialCharTok{/}\NormalTok{ prop\_complier}
\end{Highlighting}
\end{Shaded}

\begin{verbatim}
## [1] 0.1450721
\end{verbatim}

\end{document}
